\subsection{The two rows referring to the attached cyclotomic return}

\begin{result}{F8a}{Known}{The starred arithmetic centralizer assertion}
Fix $\overline{\mathbf Q}$ and write $G_{\mathbf Q}$ for its full
profinite Galois group. For every open subgroup $\Gamma'\subset G_{\mathbf Q}$,
let $\Gamma'^{\circ}$ be the kernel of its cyclotomic character. Then
\[
 C_{G_{\mathbf Q}}(\Gamma'^{\circ})=1.
\]
\passage{Cote 149, \S2, manuscript p.~22, scan 45, starred ``\`a v\'erifier !''.}
\end{result}
\begin{proof}
Let $k=\overline{\mathbf Q}^{\Gamma'}$. Galois correspondence gives the
literal equality of closed subgroups
$\Gamma'^{\circ}=G_{k(\mu_\infty)}$ in $G_{\mathbf Q}$.
Statement B1 of the supplied \texttt{CYCLOTOMIC\_KERNEL} return
\cite[Statement B1 and its proof]{PriorCyclotomic} proves
$C_{G_{\mathbf Q}}(G_{k(\mu_\infty)})=1$ for every number field $k$.
The prompt explicitly permits this citation. We checked that proof,
including the following two points essential to the proposed R3 argument.
After a finite Galois enlargement of $k$, choose a prime $p$ unramified
in that original coefficient field and split in the finite field whose
elements a proposed centralizer element moves. The value group in
$k(\mu_\infty)$, normalized by $v(p)=1$, is
$(p-1)^{-1}\mathbf Z[1/p]$. A CRT element of value one therefore has
nonzero Kummer class for a prime $\ell\nmid p(p-1)$. Equality of
stabilizers forces $\sigma(a)=c^\ell a^j$, $1\leq j<\ell$; at a
moved prime it gives $0=\ell v(c)+j$, a contradiction in that value
group. Repeating over every finite Galois number field containing $k$
forces $\sigma=1$. Thus the cited statement has exactly the subgroup,
ambient group, and quantifier required here.
\end{proof}

\begin{result}{F10a}{Proved}{The local extension and the normalizer correction}
Let $K$ be finitely generated over $\mathbf Q$, let
$U=X\setminus D$ be a smooth geometrically connected hyperbolic curve,
and let $c\in D(K)$. Choose a geometric base point and a lift of $c$
to the normalized universal covering. Its inertia $I_c$ and
decomposition group $D_c$ satisfy
\[
 D_c=N_{E_U}(I_c),\qquad
 1\longrightarrow I_c=\widehat{\mathbf Z}(1)
 \longrightarrow D_c\longrightarrow G_K\longrightarrow1,
 \qquad
 C_{E_U}(I_c)=D_c\cap p^{-1}(\ker\chi_K).
\]
A local parameter supplies a splitting. The $I_c$-conjugacy classes of
splittings form a torsor under
$H^1(K,\widehat{\mathbf Z}(1))=\varprojlim_n K^*/K^{*n}$.
\passage{Cote 149, \S3, manuscript p.~26 and continuation, scans 53--54, formulas (2)--(3).}
\end{result}
\begin{proof}
Stix \cite[Proposition 11; Corollary 12]{RootStix} applies to a
hyperbolic curve over a characteristic-zero field: distinct lifted
cusps have disjoint nontrivial inertia, the geometric normalizer of
cusp inertia is the inertia itself, and its arithmetic normalizer is
the cusp decomposition group. The tame local covers are $z^n=t$;
conjugating $z\mapsto\zeta_nz$ by an arithmetic lift sends
$\zeta_n$ to $\zeta_n^{\chi_n}$. This computes the conjugation
action and the displayed centralizer. A coefficient action fixing
compatible roots of a parameter $t$ splits the local group.
Two splittings differ by a continuous $1$-cocycle in $I_c$;
$I_c$-conjugacy changes it by a coboundary. The Kummer sequence
and Hilbert 90 give $H^1(K,\mu_n)=K^*/K^{*n}$. Passing to the limit
introduces no $\varprojlim^1H^0$ term, since the finite groups
$\mu_n(K)$ satisfy the Mittag--Leffler condition. This proves the
torsor assertion with its precise acting group.

The algebraic closure of $\mathbf Q$ in $K$ is a number field $k$,
and $K/k$ is regular. Restriction $G_K\twoheadrightarrow G_k$
therefore makes $\chi_K$ have the same nontrivial open image as
$\chi_k$. The centralizer has image $\ker\chi_K$, whereas $D_c$
surjects onto $G_K$. The word \emph{centralisateur}, retained by our
inspection of scan 54, consequently requires a mathematical correction
to \emph{normalisateur}; this is not a silent change of transcription.
\end{proof}

\begin{result}{F10b}{Known}{Embedded, disjoint cusp packets, including type $(0,3)$}
Under the hypotheses of F10a the map from local splitting classes to
$\pi_1(U_{\overline K})$-conjugacy classes of sections of $E_U\to G_K$
is injective. Its image is disjoint from the classes of points of
$U(K)$ and from every other rational cusp packet. These assertions
hold after each finite extension of $K$.
\passage{Cote 149, \S3, manuscript pp.~26--27, scans 54--55, Proposition and its margin.}
\end{result}
\begin{proof}
The exact input is Stix \cite[Theorem 14(1)--(2)]{RootStix}.
It assumes injectivity of the point-section map for smooth proper curves
of genus at least two over every finite extension of the characteristic-zero
base field, and gives the packet assertions just stated for a hyperbolic
open curve. We verify its arithmetic hypothesis. Distinct points $a,b$
of such a proper curve $C/L$ have nonzero class $[b-a]\in J_C(L)$:
a rational function with divisor $b-a$ would give a degree-one map
$C\to\mathbf P^1$. Equality of point-section classes pulls back the
finite \emph{\'etale} multiplication covers $[n]$ of $J_C$ to the
same fibre torsors; their fibre at $[a-a]=0$ is trivial. Consequently
$[b-a]\in\bigcap_n nJ_C(L)$, which is zero by the finite generation
proved and sourced in Part B. Every finite $L/K$ is again finitely
generated over $\mathbf Q$, so all the hypotheses are satisfied.
This supplies the derivation for $(0,3)$ as well as every other
hyperbolic type. The theorem is used for section classes with geometric
conjugation, exactly as in the manuscript.
\end{proof}

\begin{result}{F10c}{Refuted}{The first assertion of scan 54 in type $(0,2)$}
For $U=\mathbf G_{m,K}$ each of the two rational cusp packets equals
the entire set of section classes, and in particular contains every
point section. The first assertion of scan 54 is therefore false in
this type. For every affine characteristic-zero type with nontrivial
geometric fundamental group, the disjointness conclusion holds exactly
in the hyperbolic cases covered by F10b.
\passage{Cote 149, \S3, continuation of manuscript p.~26, scan 54.}
\end{result}
\begin{proof}
The covers $z^n=t$ identify the geometric group of $\mathbf G_m$
with $T=\widehat{\mathbf Z}(1)$. The point $1$ identifies its
arithmetic group with $T\rtimes G_K$. The local group at $0$ is
this entire semidirect product; at $\infty$ the parameter $t^{-1}$
gives the same subgroup with inertia generator inverted. Thus both
local-to-global maps on splitting classes are bijections with
$H^1(K,T)$. A point $a\in K^*$ maps to its Kummer class in that
same set. The cases $(0,0)$ and $(0,1)$ have trivial geometric group;
proper curves have no cusp packet. These computations also locate the
failure of the filling-a-hole proof in scan 55: filling one hole of
$\mathbf G_m$ gives $\mathbf A^1$, where the point-section map has
only one value.
\end{proof}
